\documentclass[11pt,a4paper]{article}

% Codificación y lenguaje
\usepackage[utf8]{inputenc}
\usepackage[T1]{fontenc}
\usepackage[spanish,es-nodecimaldot]{babel}

% Presentación
\usepackage[a4paper,margin=2.7cm]{geometry}
\usepackage{parskip}
\usepackage{fancyhdr}
\usepackage{lastpage}
\usepackage[hidelinks]{hyperref}

\pagestyle{fancy}
\fancyhf{}
\lhead{Arquitectura de Computadoras}
\rhead{TP N.\textsuperscript{o} 2 -- UART}
\cfoot{\thepage\ de \pageref{LastPage}}
\setlength{\headheight}{14pt}

\title{\textbf{Diseño e implementación de una UART}\\[0.4em]
       \large Trabajo Práctico N.\textsuperscript{o} 2}
\author{Javier Siñuka y David Trujillo}
\date{3 de septiembre de 2026}

\begin{document}

\begin{titlepage}
  \hypersetup{pageanchor=false}
  \centering
  \vspace*{1.2cm}
  {\Large Universidad \par}
  \vspace{0.3cm}
  {\large Arquitectura de Computadoras \par}
  \vspace{2.2cm}
  {\LARGE\bfseries Diseño e implementación de una UART\par}
  \vspace{0.5cm}
  {\large Trabajo Práctico N.\textsuperscript{o} 2\par}
  \vfill
  \begin{tabular}{rl}
    Estudiantes: & Javier Siñuka y David Trujillo \\
    Fecha: & 3 de septiembre de 2026
  \end{tabular}
  \vfill
\end{titlepage}

\hypersetup{pageanchor=true}
\clearpage
\pagenumbering{roman}
\tableofcontents
\clearpage
\pagenumbering{arabic}

\begin{abstract}
Este informe documentará el diseño, la implementación y la verificación del
sistema UART desarrollado para el Trabajo Práctico N.\textsuperscript{o} 2.
\end{abstract}

\section{Introducción}

\section{Objetivos}

\section{Arquitectura general}

\section{Implementación}

\section{Simulación y verificación}

\section{Conclusiones}

\end{document}
